\documentclass[12pt,a4paper]{article}
\usepackage[utf8]{inputenc}
\usepackage[spanish,es-noquoting]{babel}
\usepackage[T1]{fontenc}
\usepackage{lmodern}
\usepackage{geometry}
\geometry{top=2cm, bottom=2cm, left=2.5cm, right=2.5cm}
\usepackage{xcolor}
\usepackage{listings}
\usepackage{enumitem}
\usepackage{fancyhdr}
\usepackage{hyperref}
\usepackage{tikz}
\usetikzlibrary{arrows.meta, positioning, calc}

\definecolor{codegreen}{rgb}{0,0.6,0}
\definecolor{backcolour}{rgb}{0.95,0.95,0.92}

\lstset{
    language=C,
    backgroundcolor=\color{backcolour},
    commentstyle=\color{codegreen},
    keywordstyle=\color{blue},
    basicstyle=\ttfamily\small,
    breaklines=true,
    frame=single,
    numbers=left,
    tabsize=4,
    extendedchars=true,
    showstringspaces=false,
    inputencoding=utf8,
    literate={á}{{\'a}}1 {é}{{\'e}}1 {í}{{\'i}}1 {ó}{{\'o}}1 {ú}{{\'u}}1 {ñ}{{\~n}}1 {¡}{{\textexclamdown}}1 {¿}{{\textquestiondown}}1
}

\pagestyle{fancy}
\fancyhf{}
\lhead{Monitorías Sistemas Operativos 2026-I}
\rhead{Capítulo 3: Paso de Argumentos}
\cfoot{\thepage}

\title{\textbf{Guía de Aprendizaje: Paso de Múltiples Argumentos a Hilos}}
\author{Malak Sanchez \\ \small Monitorías de Sistemas Operativos}
\date{\today}

\begin{document}

\maketitle

\section{El Desafío}
La función \texttt{pthread\_create} solo permite pasar un único puntero (\texttt{void*}) como argumento al hilo. Si necesitamos que el hilo reciba una cadena, un entero y un flotante al mismo tiempo, no podemos hacerlo directamente.

\section{La Solución: Estructuras (\texttt{struct})}
La forma profesional y robusta de pasar múltiples datos es encapsularlos en una estructura y pasar la dirección de dicha estructura.

\begin{center}
\begin{tikzpicture}[font=\small, >=Stealth]
    % Data block
    \node[draw, fill=blue!5, minimum width=2.5cm, minimum height=1.5cm] (struct) at (0,0) {
        \begin{tabular}{c}
            \textbf{struct data} \\ \hline
            int id \\
            char* msg \\
            float value
        \end{tabular}
    };

    % Thread Node
    \node[draw, thick, fill=green!10, rounded corners=3pt, minimum width=2cm, minimum height=1cm] (thread) at (6,0) {Hilo};

    % Arrow
    \draw[->, ultra thick, violet!70!black] (struct.east) -- (thread.west) node[midway, above] {Puntero \texttt{void*}};
\end{tikzpicture}
\end{center}

\section{Recordatorio Rápido: Cómo usar \texttt{struct}}
Una estructura (\texttt{struct}) permite agrupar diferentes tipos de datos bajo un mismo nombre.

\begin{lstlisting}
// Definicion de una estructura
typedef struct {
    int edad;
    char nombre[20];
    float promedio;
} estudiante;
\end{lstlisting}

Si tiene un puntero a una estructura, use \texttt{->} en lugar de \texttt{.}

\begin{lstlisting}
estudiante* ptr = &s1;
printf("%s\n", ptr->nombre);
\end{lstlisting}

\newpage

\section{Ejemplo: Pasando Datos de un Usuario}
\begin{lstlisting}
#include <stdio.h>
#include <stdlib.h>
#include <pthread.h>

typedef struct {
    int id;
    char name[20];
    int score;
} thread_args;

void* print_user_data(void* arg) {
    thread_args* data = (thread_args*)arg; // Casting back
    printf("Thread %d: User %s has %d points.\n", data->id, data->name, data->score);
    return NULL;
}

int main() {
    pthread_t thread;
    thread_args* bundle = malloc(sizeof(thread_args));

    // Prepare data
    bundle->id = 1;
    sprintf(bundle->name, "Developer");
    bundle->score = 999;

    pthread_create(&thread, NULL, print_user_data, (void*)bundle);
    pthread_join(thread, NULL);

    free(bundle); // Clean up memory
    return 0;
}
\end{lstlisting}

\section{¡Peligro! El error del Scoping}
Nunca pase la dirección de una variable local de una función que pueda terminar antes que el hilo. Use memoria dinámica (\texttt{malloc}) para asegurar que los datos vivan lo suficiente.

\section{Ejercicios}
\begin{enumerate}
    \item Cree un programa que lance 5 hilos. A cada hilo pásele un ID único y un tiempo de espera (\texttt{sleep}) usando una estructura.
    \item Modifique el ejercicio anterior para que los hilos almacenen un resultado complejo en la misma estructura y lo lean en el \texttt{main} después del \texttt{join}.
\end{enumerate}

\end{document}
